% 编译方式: xelatex SL_ratio_proof.tex
\documentclass[UTF8]{ctexart}
\usepackage{amsmath, amssymb, amsthm}
\usepackage[margin=2.5cm]{geometry}
\usepackage[hyphens]{url}
\usepackage[hidelinks]{hyperref}
\usepackage{booktabs}
\emergencystretch 3em

\newtheorem{theorem}{定理}[section]
\newtheorem{proposition}[theorem]{命题}
\newtheorem{lemma}[theorem]{引理}
\newtheorem{remark}[theorem]{注}
\newtheorem{example}[theorem]{例}
\newtheorem{definition}[theorem]{定义}
\newtheorem{conjecture}[theorem]{猜想}

\title{相邻特征值比值 \texorpdfstring{$\lambda_{n+1}/\lambda_n$}{lambda_{n+1}/lambda_n} 的上确界定理: 证明文档}
\author{研究证明文档 (项目: BVE research)}
\date{2026-09-20 (作者修订)}

\begin{document}
\maketitle

\begin{abstract}
本文证明弦振动 Sturm-Liouville 问题 $-y'' = \lambda\rho y$ (Dirichlet 边界,
$1 \leq \rho \leq R$) 的相邻特征值比值全序列上确界定理:
\[
	\sup_{n\geq 1}\sup_{\rho} \frac{\lambda_{n+1}}{\lambda_n} = \nu(R),
\]
其中
\[
	\nu(R) = \Bigl( \frac{\arccos(-\sqrt{R}/(\sqrt{R}+1))}{\arccos(\sqrt{R}/(\sqrt{R}+1))} \Bigr)^2 .
\]
证明先在有界可测密度盒类中证明首对极值存在, 再由一阶变分、切换函数和守恒能量
推出所有极值密度的对称性及平衡块宽. 因而平衡配置的闭式计算确实给出全局极值,
不再把候选值当作上界. 随后使用配套文稿重证的 Mahar-Willner 倍指标恒等式,
结合 $\lambda_{n+1}\leq\lambda_{2n}$ 得到全序列结论.
同一首对变分论证给出 Keller 型下确界 $\mu(R)$ 的配置.
分段常数、分段连续和有界可测类的极值相同; 文中补证所需谱连续性.
证明使用经典紧自伴谱理论和 Sturm 振荡、比较理论; 本稿是修订作者的证明提交,
不是独立验收报告. 历史数值记录不作为全局最优性的证明.
固定 $n$ 的上确界猜想与全序列下确界问题均如实标注为开放, 不作为定理.
\end{abstract}

\tableofcontents

\section{记号与问题}

设 $1\leq R<\infty$, 记
\[
 \mathcal M_R=\{\rho\in L^\infty(0,1):1\leq\rho\leq R\ \text{a.e.}\},
\]
并记 $\mathcal S_R$ 为有限分段常数子类, $\mathcal{PC}_R$ 为分段连续子类.
密度按几乎处处相等识别, 跳点的单点取值不影响谱. 对 $\rho\in\mathcal M_R$, 考虑
\begin{equation}
	-y''(x) = \lambda \rho(x) y(x), \qquad y(0) = y(1) = 0 .
\end{equation}
方程按弱意义理解; 特征函数属于 $H^2(0,1)\cap H^1_0(0,1)$,
实际有 $u\in W^{2,\infty}(0,1)\subset C^1[0,1]$.
记 $\lambda_1(\rho) < \lambda_2(\rho) < \cdots$ 为特征值
(正则 SL 问题, 简单谱, 第 $m$ 个特征函数恰有 $m-1$ 个内零点).\footnote{Sturm 振荡理论: 特征值严格递增且简单, 特征函数节点数随指标递增.}
定义
\begin{equation}
	\nu(R) := \sup_{\rho} \frac{\lambda_2(\rho)}{\lambda_1(\rho)},
	\qquad
	\mu(R) := \inf_{\rho} \frac{\lambda_2(\rho)}{\lambda_1(\rho)} .
\end{equation}
此处上下确界均取遍 $\mathcal M_R$. 下述连续性说明换成另外两类时值不变.

\begin{proposition}[谱连续性、密度逼近与极值存在]\label{prop:compactness}
若 $\rho_j,\rho\in\mathcal M_R$ 且 $\rho_j\rightharpoonup^*\rho$ 于 $L^\infty$,
则对每个固定的 $k$, $\lambda_k(\rho_j)\to\lambda_k(\rho)$.
$L^1$ 收敛也有此结论. 每个固定编号的特征值比在 $\mathcal M_R$ 上达到最大和最小值,
且它在 $\mathcal S_R,\mathcal{PC}_R,\mathcal M_R$ 上的上、下确界相同.
\end{proposition}

\begin{proof}
在 $H=H^1_0(0,1)$ 上取内积 $(u,v)_H=\int_0^1u'v'$, 定义紧正自伴算子
\[
 (T_\rho u,v)_H=\int_0^1\rho uv.
\]
紧性来自 $H\hookrightarrow L^2$ 的紧嵌入, 其正特征值依次为 $1/\lambda_k(\rho)$.
$H$ 的单位球在 $C[0,1]$ 中相对紧: $|u(x)|\leq\|u'\|_2$ 且
$|u(x)-u(y)|\leq\|u'\|_2|x-y|^{1/2}$. 因此乘积族
$\{uv:\|u\|_H,\|v\|_H\leq1\}$ 在 $L^1$ 中相对紧.
弱星收敛给出对每个这样的乘积的积分收敛; $\|\rho_j-\rho\|_\infty\leq2R$
及有限网论证使收敛在乘积族上一致. 故
\[
 \|T_{\rho_j}-T_\rho\|
 =\sup_{\|u\|_H,\|v\|_H\leq1}\left|\int_0^1(\rho_j-\rho)uv\right|\longrightarrow0.
\]
紧自伴算子的极小极大原理给出每个正特征值的连续性, 取倒数得到所述结论.
若为 $L^1$ 收敛, 则直接有 $\|T_{\rho_j}-T_\rho\|\leq\|\rho_j-\rho\|_1$.

盒类 $\mathcal M_R$ 弱星闭且有界; 由弱星紧性及 $L^1$ 的可分性,
每个序列有弱星收敛子列. 又
$(k\pi)^2/R\leq\lambda_k\leq(k\pi)^2$, 故比值连续且极值存在.
最后, 在逐次二等分网格上用区间平均值逼近任意 $\rho\in\mathcal M_R$,
得到仍满足盒约束的有限阶梯密度 $\rho_j$, 且 $\|\rho_j-\rho\|_1\to0$.
此收敛可先对连续函数用一致连续性证明, 再用连续函数在 $L^1$ 中的稠密性和
区间平均算子的 $L^1$ 压缩性推出. 结合上述连续性及
$\mathcal S_R\subset\mathcal{PC}_R\subset\mathcal M_R$ 即得三个类的极值相同.
\end{proof}

\subsection{转移矩阵与 secular 方程}

在每块 $\rho \equiv c$ (宽度 $w$) 上, 相位为 $q = \sqrt{\lambda c}\,w$, 传播矩阵为
\begin{equation}
	T(q; \lambda, c) =
	\begin{pmatrix}
		\cos q & \sin q / \sqrt{\lambda c} \\
		-\sqrt{\lambda c}\,\sin q & \cos q
	\end{pmatrix},
\end{equation}
它把块左端的 $(y, y')$ 映射到块右端. 整体转移矩阵为各块矩阵按顺序之积;
Dirichlet 条件 $y(1) = 0$ 等价于总矩阵的 $(0,1)$ 元素为零 (secular 方程).
本项目所有数值结果均用该方程定位特征值, 对每个特征值局部二分加密.\footnote{数值代码见 \texttt{scripts/} 目录下的 \texttt{num\_*.py}; 例如 \texttt{num\_formula.py} 实现 $\nu, \mu$ 闭式与数值对比.}

\subsection{缩放}

$\rho \mapsto \rho/a$, $\lambda \mapsto a\lambda$ 保持方程与特征值比值不变,
故不妨设 $a = 1$, $A = R \geq 1$; 一般情形把 $R$ 换成 $A/a$.\footnote{一般 $0 < a \leq \rho \leq A$ 的比值问题与 $1 \leq \rho \leq A/a$ 等价.}

\section{首对比值的全局极值: 变分与界面论证}\label{sec:first-pair}

本节不使用倍指标引理、后文主定理或预设的平衡配置. 对称性和平衡宽度均从
任意全局极值密度的必要条件推出, 再结合极值存在性确定全局极值.

\begin{theorem}[首对全局极值的结构]\label{thm:first-pair}
设 $R>1$, $s=\sqrt R$. 在 $\mathcal M_R$ 中, $\lambda_2/\lambda_1$ 的极大密度
唯一到零测集, 为 $[1,R,1]$, 三块宽度为
\[
 \frac{s}{2s+1},\quad\frac1{2s+1},\quad\frac{s}{2s+1}.
\]
极小密度唯一到零测集, 为 $[R,1,R]$, 三块宽度为
\[
 \frac1{s+2},\quad\frac{s}{s+2},\quad\frac1{s+2}.
\]
因此任意容许密度的首对比值不超过前一配置的比值, 不低于后一配置的比值.
\end{theorem}

\begin{proof}
\textbf{1. 一阶变分与切换规则.}
命题~\ref{prop:compactness} 保证最大、最小密度均存在. 固定其中任一个 $\rho$,
取实特征函数 $u_1,u_2$, 使 $\int_0^1\rho u_i^2=1$, $u_1>0$ 于 $(0,1)$,
且 $u_2$ 在其唯一内零点左侧为正. 对任意 $\sigma\in\mathcal M_R$,
令 $h=\sigma-\rho$, $\rho_t=\rho+th$, $0\leq t\leq1$.
简单特征值的参数微分给出
\begin{equation}\label{eq:first-variation}
 \dot\lambda_i=-\lambda_i\int_0^1hu_i^2,
 \qquad \left.\frac{d}{dt}\log\frac{\lambda_2(\rho_t)}{\lambda_1(\rho_t)}\right|_{t=0}
 =\int_0^1hF,\qquad F=u_1^2-u_2^2.
\end{equation}
这与 Keller~\cite[pp.~486--487, (3.4)--(3.6)]{keller} 的变分公式一致;
下面给出在本稿可测盒类中的论证, 不把原始密度范围直接扩大.
这里的参数微分对 $L^\infty$ 密度成立: 射击解满足 Volterra 积分方程,
对 $t,\lambda$ 可微; 在 Dirichlet 根处, 令 $v=\partial_\lambda y$, 有
$(y'v-yv')'=\rho y^2$, 故
$y'(1)v(1)=\int_0^1\rho y^2>0$, 特征方程的根简单.
隐函数定理给出可微特征对. 将其方程微分、乘 $u_i$ 并分部积分,
边界项为零, 即得~\eqref{eq:first-variation}.

最大值处 $\int(\sigma-\rho)F\leq0$ 对所有 $\sigma$ 成立, 故
\begin{equation}\label{eq:switch-rule}
 \rho=R\ \text{于 }\{F>0\},\qquad \rho=1\ \text{于 }\{F<0\}\quad\text{a.e.}
\end{equation}
例如在 $\{F>0\}$ 上把 $\sigma$ 改为 $R$ 即证明第一条.
最小值处变分不等号相反, 因而两条规则中的 $1,R$ 交换.

\textbf{2. 切换函数至多有两个内零点.}
令 $\xi$ 为 $u_2$ 的唯一内零点, 并置
\[
 W=u_1u_2'-u_1'u_2,\qquad r=u_2/u_1.
\]
有 $W'=-(\lambda_2-\lambda_1)\rho u_1u_2$ a.e., $W(0)=W(1)=0$.
在 $(0,\xi)$ 上从左端积分得 $W<0$; 在 $(\xi,1)$ 上从右端反向积分仍得 $W<0$.
所以 $r'=W/u_1^2<0$ 于整个 $(0,1)$.
而 $F=u_1^2(1-r^2)$, 故 $F=0$ 至多发生在 $r=1,-1$ 两点,
特别地没有正测度的奇异切换集.
该商函数机制亦见 Mahar--Willner~\cite[pp.~526--527, Lemma 4,
(5.21)--(5.23)]{mw}; 这里的双端积分直接给出整个开区间的严格单调性.

\textbf{3. 能量恒等为零与端点斜率.}
在最大值情形定义 $g(t)=Rt$ ($t\geq0$), $g(t)=t$ ($t\leq0$);
最小值情形定义 $g(t)=t$ ($t\geq0$), $g(t)=Rt$ ($t\leq0$).
切换规则给出 $g(F)=\rho F$ a.e. 定义连续的能量代表
\begin{equation}\label{eq:zero-energy}
 E=\frac{u_1'^2}{\lambda_1}-\frac{u_2'^2}{\lambda_2}+g(F).
\end{equation}
$u_i'\in W^{1,\infty}$, $g$ 为 Lipschitz, 故 $E$ 绝对连续.
除至多两个内零点外, $[g(F)]'=\rho F'$ a.e., 从微分方程得
\[
 E'=-2\rho u_1u_1'+2\rho u_2u_2'+\rho F'=0\quad\text{a.e.}
\]
这也解释了密度跳跃处为何没有遗漏的界面项: $g(F)$ 连续, 且切换时 $F=0$.
因此 $E$ 为常数; 用归一化和分部积分,
\[
 \int_0^1E=\frac{\int u_1'^2}{\lambda_1}
 -\frac{\int u_2'^2}{\lambda_2}+\int\rho(u_1^2-u_2^2)=1-1+1-1=0.
\]
故 $E\equiv0$. 令 $q=\sqrt{\lambda_2/\lambda_1}>1$.
端点处 $F=0$, 非零端点导数及两函数的符号给出
\begin{equation}\label{eq:endpoint-ratios}
 \frac{u_2'(0)}{u_1'(0)}=q,\qquad
 \frac{u_2'(1)}{u_1'(1)}=-q.
\end{equation}
于是 $r(0+)=q>1$, $r(1-)=-q<-1$.
结合严格递减性, 恰有 $0<\alpha<\xi<\beta<1$ 使 $r(\alpha)=1$,
$r(\beta)=-1$; $F<0$ 在两端区间, $F>0$ 在中间区间.
最大密度因此为 $[1,R,1]$, 最小密度为 $[R,1,R]$, 但此时尚未假设对称或平衡宽度.

\textbf{4. 由外侧匹配推出对称性.}
统一记外侧密度为 $d_o$, 中间密度为 $d_m$, 两者在最大情形为 $(1,R)$,
在最小情形为 $(R,1)$, 并令 $p_i=\sqrt{\lambda_i d_o}$.
由~\eqref{eq:endpoint-ratios}, 在左外块有同一个常数 $A_0>0$ 使
\[
 u_1(x)=A_0\sin(p_1x),\qquad u_2(x)=A_0\sin(p_2x).
\]
在 $x=\alpha$, 两函数相等且为正. 因左外块内两者均无零点,
$0<p_1\alpha<p_2\alpha<\pi$.
所以 $\sin(p_1\alpha)=\sin(p_2\alpha)$ 只能给出
\[
 (p_1+p_2)\alpha=\pi.
\]
从右端对 $u_1(1-x)$ 与 $-u_2(1-x)$ 作同样论证得
$(p_1+p_2)(1-\beta)=\pi$. 因而 $\beta=1-\alpha$, 密度关于中点对称.
简单谱和反射不变性给出 $u_1$ 关于中点为偶函数, $u_2$ 为奇函数.
后一结论也可由其恰有一个简单内零点看出: 偶函数的非中心零点成对,
而中心零点会同时有零导数. 故 $\xi=1/2$.

\textbf{5. 由内侧匹配推出平衡宽度, 不遗漏三角分支.}
令 $h=1/2-\alpha>0$, $X=p_1\alpha=\pi/(1+q)\in(0,\pi/2)$,
$Y=\sqrt{\lambda_1d_m}\,h$, $Z=\sqrt{\lambda_2d_m}\,h$.
中间块的偶、奇解分别是余弦和正弦; $u_1>0$ 给出 $0<Y<\pi/2$,
$u_2>0$ 于 $(\alpha,1/2)$ 给出 $0<Z<\pi$.
在 $\alpha$ 处连续匹配对数导数, 并用 $p_2\alpha=\pi-X$, 得
\begin{equation}\label{eq:interface-phases}
 \cot X=\sqrt{d_m/d_o}\tan Y
       =\sqrt{d_m/d_o}\cot Z.
\end{equation}
因为 $\cot X>0$, 必有 $Z<\pi/2$. 因而 $\tan Y=\cot Z$ 唯一地给出
$Y+Z=\pi/2$. 与 $(p_1+p_2)\alpha=\pi$ 比较得
\[
 \sqrt{d_o}\,\alpha=2\sqrt{d_m}\,h
 =\sqrt{d_m}(1-2\alpha),\qquad
 \alpha=\frac{\sqrt{d_m}}{2\sqrt{d_m}+\sqrt{d_o}}.
\]
代入两组 $(d_o,d_m)$ 就得到定理的两种块宽.

最后, 上述推导适用于每个全局极大或极小密度, 而这两种极值已由紧性保证存在.
故它们分别只能是所列的唯一配置. 这一步给出对全部容许密度的上、下界,
并非仅仅验证某个平衡配置满足驻点条件.
\end{proof}

\begin{remark}[退化参数]
$R=1$ 时唯一密度为 $1$ a.e., $\lambda_j=j^2\pi^2$, 首对上、下确界均为 $4$.
上述三块公式在 $R=1$ 只表示同一个常密度, 不再有实际跳点.
\end{remark}

\section{主要定理}

\begin{theorem}[相邻比值上确界]
对 $R \geq 1$,
\begin{equation}
	\boxed{\; \sup_{n \geq 1}\, \sup_{1 \leq \rho \leq R} \frac{\lambda_{n+1}(\rho)}{\lambda_n(\rho)}
	= \nu(R) = \left( \frac{\arccos\!\bigl(-\sqrt{R}/(\sqrt{R}+1)\bigr)}
	{\arccos\!\bigl(\sqrt{R}/(\sqrt{R}+1)\bigr)} \right)^2 . \;}
\end{equation}
上确界在 $n = 1$ 处由 $\rho_* = [1, R, 1]$ 达到, 即
\[
	\rho_*(x) =
	\begin{cases}
		1, & x \in [0, st] \cup [st + t, 1], \\
		R, & x \in (st, st + t),
	\end{cases}
	\qquad s = \sqrt{R}, \quad t = \frac{1}{2\sqrt{R}+1}.
\]
\end{theorem}

\begin{proof}
\textbf{第一步: 首对全局上界.} 定理~\ref{thm:first-pair} 证明任意容许密度的
$\lambda_2/\lambda_1$ 不超过 $\rho_*$ 的比值. 第~\ref{sec:balanced} 节计算该比值为
所列闭式, 因而这次确有 $\nu(R)=\lambda_2(\rho_*)/\lambda_1(\rho_*)$.
$R=1$ 的结论直接由常密度谱给出.

\textbf{第二步: 平凡不等式.} 对固定的 $\rho$, 特征值序列严格递增, 而 $n + 1 \leq 2n$
对一切 $n \geq 1$ 成立, 故 $\lambda_{n+1}(\rho) \leq \lambda_{2n}(\rho)$, 从而
\begin{equation}
	\frac{\lambda_{n+1}(\rho)}{\lambda_n(\rho)} \leq \frac{\lambda_{2n}(\rho)}{\lambda_n(\rho)} .
\end{equation}

\textbf{第三步: 倍指标恒等式 (Mahar-Willner).} 配套文稿
\texttt{SL\_mw\_lemma\_reproof.tex} 的 max/min doubling theorems 重证了
Mahar-Willner 型恒等式: 对每个 $n\geq1$,
\begin{equation}
	\sup_{\rho} \frac{\lambda_{2n}(\rho)}{\lambda_n(\rho)} = \nu(R),
	\qquad
	\inf_{\rho} \frac{\lambda_{2n}(\rho)}{\lambda_n(\rho)} = \mu(R).
\end{equation}
配套证明以有限阶梯密度为起点, 使用有符号端点斜率拼接、Wronskian 零点运动和
穷尽共同零点情形的截断归纳, 再经 $L^1$ 逼近推广到 $\mathcal M_R$.
配套文稿使用 $[a,1]$ 归一化; 本处取 $a=1/R$ 并将密度除以 $R$,
所以该稿的 $\nu(a),\mu(a)$ 分别对应此处的 $\nu(R),\mu(R)$.
其中 $\nu,\mu$ 只按首对上、下确界定义, 不使用它们的闭式或本定理.
故该引用与第一步的首对变分论证没有循环.

合并上述不等式得
\begin{equation}
	\sup_{n \geq 1} \sup_{\rho} \frac{\lambda_{n+1}(\rho)}{\lambda_n(\rho)}
	\leq \sup_{n \geq 1} \sup_{\rho} \frac{\lambda_{2n}(\rho)}{\lambda_n(\rho)}
	= \nu(R).
\end{equation}

第一步已证明取 $n=1$ 和 $\rho_*$ 达到该上界.
\end{proof}

\begin{remark}[适用范围与依赖]
证明链为: 本稿首对变分结构定理 $\to$ 平衡配置闭式;
配套文稿的倍指标恒等式 $\to$ 全序列上界; 两者在本定理汇合.
经典紧自伴谱理论、弱星紧性和 Sturm 理论是背景依赖, 特定极值定理不作为黑箱前提.
这是作者给出的证明链, 是否验收由独立复核决定.
定理不声称 $n \geq 2$ 时 $\lambda_{n+1}/\lambda_n$ 的极值, 那里只有数值猜想的
$\Lambda_n^{\sup}(R)$ (第~\ref{sec:open} 节), 本稿未解决. 特别地, 本定理不蕴含任何
关于 $\inf \lambda_{n+1}/\lambda_n$ 的结论.
\end{remark}

\section{闭式推导: 平衡相位方法}\label{sec:balanced}

本节只计算指定配置的比值, 分别记为 $C_+(R),C_-(R)$.
定理~\ref{thm:first-pair} 已证明极值配置必有这些平衡块宽; 结合该定理才能识别
$\nu(R)=C_+(R)$ 和 $\mu(R)=C_-(R)$.
计算本身使用相同块相位 $p=\sqrt{\lambda\rho}\cdot\text{块宽}$,
把 secular 方程化为纯三角方程.

\subsection{上确界配置 \texorpdfstring{$[1,R,1]$}{[1,R,1]}}

设块宽 $w_1, w_2, w_3$ (权重 $1, R, 1$). 平衡条件 $q_1 = q_2 = q_3 =: p$ 即
$\sqrt{\lambda}\,w_1 = \sqrt{\lambda R}\,w_2 = \sqrt{\lambda}\,w_3$, 故 $w_1 = w_3 = s w_2$,
$s = \sqrt{R}$. 总宽 $(2s+1)w_2 = 1$ 给出 $w_2 = t = 1/(2s+1)$, $w_1 = w_3 = st$:
恰为 $\rho_*$. 此时 $\lambda = (p/(st))^2$.

\begin{lemma}[平衡相位 secular 方程, 上确界配置]
对 $[1,R,1]$ (块宽 $st, t, st$, 权重 $1, R, 1$), 平衡相位 $p$ 下 Dirichlet 条件等价于
\begin{equation}
	\sin p \cdot \Bigl( (2s+1)\cos^2 p - s^2 \sin^2 p \Bigr) = 0 .
 \label{eq:balanced-secular}
\end{equation}
\end{lemma}

\begin{proof}
直接乘转移矩阵. 记 $a = \cos p$, $b = \sin p$, $\ell = \sqrt{\lambda} = p/(st)$.
各块矩阵为 $M_1 = M_3 = \begin{pmatrix} a & bst/p \\ -bp/(st) & a \end{pmatrix}$,
$M_2 = \begin{pmatrix} a & bt/p \\ -bp/t & a \end{pmatrix}$.
先乘 $M_2 M_1$: 左上元 $a^2 - b^2/s$, 右上元 $ab\,t(s+1)/p$,
左下元 $-ab\,p(s+1)/(st)$, 右下元 $a^2 - b^2 s$. 再乘 $M_1 (M_2 M_1)$ 并取 $(0,1)$ 元:
\begin{equation}
	y(1) = \frac{b t}{p}\Bigl( (2s+1)a^2 - s^2 b^2 \Bigr),
\end{equation}
即~\eqref{eq:balanced-secular}; 这是对 $p>0$ 的代数恒等式.
\end{proof}

解~\eqref{eq:balanced-secular} 的非平凡根: $(2s+1)\cos^2 p = s^2 \sin^2 p$, 即
\begin{equation}
	\cos^2 p = \frac{s^2}{(s+1)^2}
	\quad\Longleftrightarrow\quad
	\cos p = \pm \frac{s}{s+1} .
\end{equation}
于是前两个特征值对应相位
\begin{equation}
	p_1 = \theta := \arccos\!\Bigl( \frac{s}{s+1} \Bigr) \in (0, \pi/2),
	\qquad
	p_2 = \pi - \theta ,
\end{equation}
(在 $(0, \pi)$ 内 $\cos^2 p = (s/(s+1))^2$ 恰有两个根 $\theta, \pi-\theta$;
$\sin p = 0$ 的根 $p = \pi$ 对应更高特征值). 因此
\begin{equation}
	\lambda_1 = \Bigl( \frac{(2s+1)\theta}{s} \Bigr)^2,
	\qquad
	\lambda_2 = \Bigl( \frac{(2s+1)(\pi-\theta)}{s} \Bigr)^2,
\end{equation}
\begin{equation}
 C_+(R)=\frac{\lambda_2(\rho_*)}{\lambda_1(\rho_*)}
 =\Bigl(\frac{\pi-\theta}{\theta}\Bigr)^2
 =\Bigl(\frac{\arccos(-s/(s+1))}{\arccos(s/(s+1))}\Bigr)^2.
\end{equation}
这同时验证了恒等式 $\sqrt{\lambda_1}\,st = \theta$ 与 $\sqrt{\lambda_1} = (2s+1)\theta/s$.
至此再应用定理~\ref{thm:first-pair}, 得 $\nu(R)=C_+(R)$;
若没有该结构定理, 以上计算只能给出 $C_+(R)\leq\nu(R)$.

\subsection{Keller 下确界配置 \texorpdfstring{$[R,1,R]$}{[R,1,R]}}

对权重 $R, 1, R$ 与块宽 $c, sc, c$, 平衡条件给出 $c = 1/(s+2)$,
$\lambda = (p/(sc))^2$. 记 $\ell=\sqrt\lambda$, 同理计算转移矩阵, 得
\begin{equation}
	y(1) = \frac{\sin p}{\ell s^2}\Bigl( s(s+2)\cos^2 p - \sin^2 p \Bigr) = 0
	\quad\Longleftrightarrow\quad
	\tan^2 p = s(s+2) = R + 2\sqrt{R} .
\end{equation}
取 $\varphi := \arccos(1/(s+1))$, 则 $\tan^2\varphi = (s+1)^2 - 1 = s^2 + 2s = s(s+2)$,
故前两个特征值对应 $p = \varphi$ 与 $p = \pi - \varphi$:
\begin{equation}
 C_-(R)=\frac{\lambda_2}{\lambda_1}=\Bigl(\frac{\pi-\varphi}{\varphi}\Bigr)^2
 =\Bigl(\frac{\arccos(-1/(s+1))}{\arccos(1/(s+1))}\Bigr)^2,
\end{equation}
同样, 必须结合定理~\ref{thm:first-pair} 才得到 $\mu(R)=C_-(R)$.
同时 $\sqrt{\lambda_1}=\varphi(s+2)/s$.

\begin{remark}[与数值的对照]
转移矩阵数值对 $R = 1.5, 2, 3, 4, 10, 100$ 复现 $\nu(R), \mu(R)$ 闭式至 8 位小数;
角度恒等式 $\sqrt{\lambda_1}\,st = \theta$, $\sqrt{\lambda_2}\,st = \pi-\theta$,
$\sqrt{\lambda_1}\,sc = \varphi$, $\sqrt{\lambda_2}\,sc = \pi-\varphi$ 对 $R = 2, 4, 10$
验证到 $10^{-15}$ 量级. 注意交接摘要中曾出现 $\sqrt{\lambda_1} = 2(s+1)\varphi/s$ 的笔误
(仅 $R = 4$ 时与正确式 $(s+2)\varphi/s$ 巧合), 本文一律使用已核验的恒等式.
\end{remark}

\section{首对下确界与文献来源边界}

\begin{theorem}[Keller, 闭式形式]
对本稿 Dirichlet 问题,
\[
 \inf_{\rho\in\mathcal M_R}\lambda_2/\lambda_1=\mu(R)
 =\left(\frac{\arccos(-1/(\sqrt R+1))}{\arccos(1/(\sqrt R+1))}\right)^2,
\]
极小化配置为 $[R, 1, R]$ (块宽 $c, sc, c$, $c = 1/(\sqrt{R}+2)$).
\end{theorem}

\begin{proof}
$R>1$ 时由定理~\ref{thm:first-pair} 和第~\ref{sec:balanced} 节的 $C_-(R)$ 计算得到.
$R=1$ 时直接取常密度.
\end{proof}

本轮已取得本地一手全文文件 \texttt{papers/mw1976.pdf} 和
\texttt{papers/keller1976.pdf} (路径相对项目根), 按下列定位阅读正文并核对 PDF 原页;
同名 \texttt{.txt} 只用于检索定位, 不作为公式无误的保证.
Mahar--Willner~\cite[pp.~517--518, Sections 1--3]{mw} 的原始范围是
$[-1/2,1/2]$ 上跳点数有界的分段连续密度 $0<a\leq\rho\leq1$.
p.~518 的 Theorems 1--2 给出极值密度至多两跳且对称;
紧接其后的文字明确未断言跳点唯一性.
最大配置的界面与归一化方程位于 p.~519 的 (4.1)--(4.5).
倍指标传递为 p.~519 的 Theorem 3, 证明见 pp.~521--524 的 Lemmas 1--2
及 p.~529. 原文胞元公式 (5.3) 位于 p.~522, 第一、第二特征函数的系数分别为
$(-s)^j$ 和 $t^j$, 不能把第二个符号改成 $(-t)^j$ 且仍要求 $t>0$.

Keller~\cite[p.~485, (1.1)--(2.1)]{keller} 的原始范围是跳点数有界的分段连续密度
$0<a\leq\rho\leq A$. 一阶变分与切换条件在 pp.~486--487 的 (3.1)--(3.9),
p.~488 明确把两跳对称结构的证明归于 Mahar--Willner;
极小配置与界面条件见 pp.~488--489 的 (4.1)--(4.8).
不能把这个结构归属省略后声称 Keller 独自给出了所有结构证明.

本稿首对全局上界和唯一性的直接证明仍为第~\ref{sec:first-pair} 节,
不只引用文献结构定理后代入一个平衡候选. 可测类的紧性与谱连续性由
命题~\ref{prop:compactness} 承担, 不假定原文的分段连续类本身具有所需弱星紧性.
原文定位用于来源与方法核对, 不主张本稿闭式或推导具有新颖性.
本次没有对两篇原文全部论证作独立验收.

\section{周期延拓构造与数值复现}

\subsection{胞元与合并构造}

设胞元为 $\phi_0 = [1, R, 1]$ 限制在 $[-1/2, 1/2]$ 上的缩放 (块宽 $st, t, st$,
总宽 $1$). 在 $x\in[j/n,(j+1)/n]$, $j=0,\ldots,n-1$, 置
$\phi_n(x)=\phi_0(nx-j-1/2)$, 得到 $[0,1]$ 上 $n$ 个胞元的重复.
相邻胞元的两个 $1$-块可合并成一个宽度 $2st/n$ 的块, 便于只记录实际跳点.
合并只是同一密度的不同表示, 不改变传递矩阵或谱, 也不是胞元恒等式成立的条件.

合并后的配置: 共 $2n+1$ 块, 模式 $[1, R, 1, R, \dots, 1]$,
其中 $n$ 个 $R$-块宽 $t/n$, 两个端 $1$-块宽 $st/n$, 其余 $n-1$ 个内部 $1$-块宽 $2st/n$.

\subsection{平方律}

对上述配置, 数值验证 (R=4):
\begin{equation}
	\lambda_n = n^2 \lambda_1^{(\text{cell})}, \qquad
	\lambda_{2n} = n^2 \lambda_2^{(\text{cell})}, \qquad
	\frac{\lambda_{2n}}{\lambda_n} = \frac{\lambda_2}{\lambda_1} = \nu(4) = 7.48153339,
\end{equation}
原稿记录对 $n=1,\ldots,6$ 的数值误差在 $10^{-8}$ 内.
本轮的解析依据是配套文稿的有符号胞元拼接, 数值记录不代替证明.
这里的平方因子是重复数 $n^2$; 不可由此推出 $\lambda_{kn}=k^2\lambda_n$,
因为一般胞元并无 $\lambda_k^{(\mathrm{cell})}=k^2\lambda_1^{(\mathrm{cell})}$.

\section{开放问题 (不作为定理)}\label{sec:open}

\begin{conjecture}[固定 $n$ 上确界]
对每个 $n \geq 1$, $\sup_{\rho} \lambda_{n+1}/\lambda_n$ 由交替 bang-bang 配置
$[1, R, 1, R, \dots, 1]$ ($2n+1$ 块, $w_1/w_R = \sqrt{R}$, $w_R = 1/((n+1)\sqrt{R}+n)$) 达到.
\end{conjecture}

该猜想未证明. 数值支持 (R=4): $c_1 = 7.48153$, $c_2 = 4.28466$, $c_3 = 3.45388$,
$c_4 = 3.09118$, $c_5 = 2.89443$, $c_6 = 2.77394$, $c_8 = 2.63833$, $c_{10} = 2.56698$,
$c_{12} = 2.52461$, $c_{16} = 2.47873$, $c_{20} = 2.45566$, $c_{24} = 2.44243$;
Keller 变分条件 (跳点处 $|y_n| = |y_{n+1}|$ 型匹配) 对 $n = 1,\dots,8$ 符号级验证 (1e-11).
能带极限 $c_{\infty}(R) = ((\pi-\varphi_1)/\varphi_1)^2$,
$\varphi_1 = \arccos((\sqrt{R}-1)/(\sqrt{R}+1))$; $c_{\infty}(4) = \mu(4) = 2.40916855$
系 $R = 4$ 巧合 ($c_{\infty}(2) = 1.55403629 \neq \mu(2)$).

\begin{remark}[下确界开放]
$\inf_{n\geq 1}\inf_{\rho}\lambda_{n+1}/\lambda_n$ 是开放问题.
已证: $n=1$ 时为 $\mu(R)$ (Keller); 一切 $n$ 下比值 $> 1$ (简单谱).
数值: $R=4$ 时交替族给出 $\lambda_5/\lambda_4 = 1.08380983 < \mu(4)$, 且随 $n$ 递减.
能否趋于 $1$ 未知.
\end{remark}

\section{历史数值记录}

以下保留原稿数值表和样例. 本轮局部核对及其范围另记于外部作者报告
\texttt{ratio-author.md}; 不把未重新运行的历史脚本记录当作本轮独立验收.

\begin{center}
\begin{tabular}{cccccc}
	\toprule
	$R$ & $\nu(R)$ 闭式 & $\nu(R)$ 数值 & $\mu(R)$ 闭式 & $\mu(R)$ 数值 \\
	\midrule
	1.5 & 4.75382078 & 4.75382078 & 3.40068243 & 3.40068243 \\
	2.0 & 5.40388487 & 5.40388487 & 3.05139810 & 3.05139810 \\
	3.0 & 6.51973823 & 6.51973823 & 2.64587239 & 2.64587239 \\
	4.0 & 7.48153339 & 7.48153339 & 2.40916855 & 2.40916855 \\
	10.0 & 11.82037283 & 11.82037283 & 1.86419353 & 1.86419353 \\
	100.0 & 39.83043629 & 39.83043629 & 1.26121838 & 1.26121838 \\
	\bottomrule
\end{tabular}
\end{center}

反例 (R=4, 交替 $[R,1,R,1,\dots,R]$, 宽度 $w_R = 1/(3n+1)$, $w_1 = 2w_R$):
$n=2$: $\lambda_3/\lambda_2 = 1.4242433972$; $n=3$: $\lambda_4/\lambda_3 = 1.1791885971$
($\lambda_3 = 56.70885901$, $\lambda_4 = 66.87043990$);
$n=4$: $\lambda_5/\lambda_4 = 1.0838098314$
($\lambda_4 = 100.09539937$, $\lambda_5 = 108.48437791$). 均小于 $\mu(4) = 2.40916855$.

\begin{thebibliography}{99}

\bibitem{mw} T. J. Mahar, B. E. Willner, \emph{An extremal eigenvalue problem},
	Comm. Pure Appl. Math. 29 (1976), 517--529.
	DOI: \url{https://doi.org/10.1002/cpa.3160290505}

\bibitem{keller} J. B. Keller, \emph{The minimum ratio of two eigenvalues},
	SIAM J. Appl. Math. 31 (1976), no.~3, 485--491.
	DOI: \url{https://doi.org/10.1137/0131042}

\end{thebibliography}

\section{涉及到的数学知识}

\begin{description}
	\item[转移矩阵与 secular 方程] 每块上的解由 $2\times 2$ 传播矩阵推进,
		Dirichlet 条件化为总矩阵元的零点方程; 全部数值基于此.
	\item[Sturm 振荡理论] 第 $m$ 个特征函数恰有 $m-1$ 个内零点; 简单谱.
	\item[Pr\"ufer 相位] 极坐标变换把特征值问题化为相位单调方程, 是平衡相位方法的背景.
	\item[谱单调性] $n+1 \leq 2n$ 推出 $\lambda_{n+1} \leq \lambda_{2n}$, 平凡但关键.
	\item[Keller 变分条件] 极值配置跳点处的匹配条件 (权重归一化 $\int \rho y^2 = 1$).
	\item[bang-bang 原理] 逐点夹逼类 $[1,R]$ 的极值配置逐点取端点值.
	\item[周期延拓与平方律] $\phi_n$ 周期化后 $\lambda_n = n^2\lambda_1^{(\text{cell})}$,
		$\lambda_{2n} = n^2\lambda_2^{(\text{cell})}$; 胞界用有符号端点斜率保证 $C^1$ 拼接.
	\item[Bloch 能带] 周期介质谱为能带, 带内比值 $\to 1$, 带边给出 $c_{\infty}(R)$.
	\item[零点截断归纳] MW Lemma 2 的核心技巧, 用于传递极值常数.
	\item[Liouville 变换与 Feynman-Hellmann] 一般 SL 标准化与一阶变分工具, 供后续推广.
\end{description}

\end{document}
